\section{The finite Paley--Wiener projection and its Gram operator}
\label{sec:global-gram}

The finite-level orthogonal projector is a concrete additional
operator.  The same range can be parametrized by a local theta
map, by a Hardy-space inner multiplier, or by frequency data,
but these parametrizations are not automatically isometries for
the specified global norm.  This section gives the actual Gram
factor, with its domain and inverse, and an explicit finite-level
theta Gram kernel.  It does not infer an unbounded finite-level
inverse from the uncut obstruction MA-MO.

\paragraph{MA-GP-004 (the exact bounded Gram correction on a half-line).}
Fix $\alpha>0$, $s\ge0$ and $c\in\mathbb R$.  On $[c,\infty)$
put
\[
 \begin{aligned}
 w_+(x)&=e^{2\alpha x}(1+x^2)^{s/2},&
 w(x)&=e^{2\alpha|x|}(1+x^2)^{s/2},\\
 \mathcal K_+&=L^2([c,\infty),w_+(x)\,dx),&
 R(x)&=w(x)/w_+(x)=e^{2\alpha|x|-2\alpha x}.
 \end{aligned}
 \tag{MA-GP.40}
\]
Let $\mathcal M\subset\mathcal K_+$ be the closed finite-level
subspace supplied by the actual smooth PW range, and let
$i:\mathcal M\hookrightarrow\mathcal K_+$ be its isometric
inclusion for the $w_+$ norm.  The two norms on this half-line
are equivalent, since
\[
            1\le R(x)\le C_c:=e^{4\alpha\max(-c,0)}.
 \tag{MA-GP.41}
\]
Consequently either norm gives the same range closure.  The
positive operator
\[
              G_c=i^*M_Ri:\mathcal M\to\mathcal M
 \tag{MA-GP.42}
\]
is boundedly invertible, with $1\le G_c\le C_c$ and
$C_c^{-1}\le G_c^{-1}\le1$.  The orthogonal projection for
the \emph{actual} $w$ norm is exactly
\[
                  Q_c^{(w)}=iG_c^{-1}i^*M_R.
 \tag{MA-GP.43}
\]
On the whole global space first restrict the input to
$[c,\infty)$ and extend the output by zero.  Formula
(MA-GP.43) then gives the actual global finite-level projector.

\begin{proof}
The weight ratio is (MA-GP.41), including for every polynomial
weight $s$.  For $u\in\mathcal M$,
$\langle G_cu,u\rangle_+=\|iu\|_w^2$, which lies between
$\|u\|_+^2$ and $C_c\|u\|_+^2$.  The spectral theorem
therefore gives the asserted bounded inverse.
For $h\in\mathcal K_+$ the vector in (MA-GP.43) belongs
to $\mathcal M$, and for every $v\in\mathcal M$,
\[
 \langle h-iG_c^{-1}i^*M_Rh,iv\rangle_w
   =\langle i^*M_Rh-G_cG_c^{-1}i^*M_Rh,v\rangle_+=0.
 \tag{MA-GP.44}
\]
This is the defining property of the weighted orthogonal
projection.  The two support half-lines are orthogonal in
the ambient weighted space, so restriction and zero extension
prove the final assertion.
\end{proof}

For $s=0$ the actual Hardy identification of MA-GP-003,
in its stated range $0<\alpha<1/2$ and $c=-\log Y$,
makes this formula explicit.  Write
\[
 (U_cf)(r)=e^{\alpha(c+r)}f(c+r),\quad r\ge0,
 \qquad
 (\mathcal Lv)(z)=\int_0^\infty e^{-zr}v(r)\,dr.
 \tag{MA-GP.45}
\]
The boundary norm $(2\pi)^{-1}\int|F(it)|^2\,dt$, as in
MA-GP-003, makes $\mathcal L$ unitary onto
$H^2(\operatorname{Re}z>0)$.  If $M_{B_\alpha}$ is the
isometric inner multiplier describing the actual finite PW
range, define
\[
 \mathcal R_c=\mathcal L U_cM_RU_c^{-1}\mathcal L^{-1},
 \qquad
 \mathcal G_c=M_{B_\alpha}^*\mathcal R_cM_{B_\alpha}.
 \tag{MA-GP.46}
\]
In this Hardy picture the projector is
\[
 M_{B_\alpha}\mathcal G_c^{-1}
                   M_{B_\alpha}^*\mathcal R_c.
 \tag{MA-GP.47}
\]
The zeros of $B_\alpha$ are those of MA-GP-003, namely
$z=\alpha+1/2-\rho$ in the right half-plane, with full
multiplicity.  Thus left-exterior global zeros can enter this
finite-level inner factor even when they do not survive in
the completed full-strip quotient of MA-WP.

There is an explicit kernel for the weight correction.
If $c<0$ and $L=-c$, then $\mathcal R_c=1+\mathcal K_c$,
where on Hardy boundary values
\[
 (\mathcal K_cF)(t)=\frac1{2\pi}\int_{\mathbb R}
                       \kappa_c(t-u)F(u)\,du,
 \qquad
 \kappa_c(v)=\int_0^L
            \bigl(e^{4\alpha(L-r)}-1\bigr)e^{-ivr}\,dr.
 \tag{MA-GP.48}
\]
The convolution is interpreted in the bounded $L^2$ sense,
as the Fourier transform of multiplication by the displayed
bounded compactly supported function.  Hence
\[
       \mathcal G_c=1+M_{B_\alpha}^*\mathcal K_cM_{B_\alpha},
       \qquad 1\le\mathcal G_c\le Y^{4\alpha}
                    \quad(Y>1).
 \tag{MA-GP.49}
\]
For $Y\le1$ the ratio $R$ is one and this particular weight
Gram is the identity.  The Hardy inner factor remains part
of the projection.  For $s>0$, (MA-GP.43) remains valid, but
no replacement of the actual Hardy--Sobolev range by the
$s=0$ inner-multiplier model is asserted.

\paragraph{The Gram factor changes the projection.}
Suppose $c<0$, $s=0$, and the inner factor has a zero and
has no positive delay, as in MA-GP-003.  Then the weighted
projection (MA-GP.47) differs from the ordinary Hardy
projection $M_{B_\alpha}M_{B_\alpha}^*$.

\begin{proof}
If the projections agreed, their range and ordinary orthogonal
complement would be orthogonal also for the $R$ inner product.
Thus multiplication by $R(c+r)$ would commute with the
ordinary projection onto the inverse-Laplace image of
$B_\alpha H^2$.  Its complement would reduce that multiplier.
On $0<r<L=-c$, $R(c+r)=e^{4\alpha(L-r)}$ is strictly
monotone.  Its spectral projections therefore include
multiplication by $\mathbf1_{(a,b)}$ for every
$0<a<b<L$.

If $B_\alpha(\mu)=0$ with $\operatorname{Re}\mu>0$,
the Hardy reproducing kernel at $\mu$ belongs to the
orthogonal complement of $B_\alpha H^2$.  Its inverse
Laplace transform is a nonzero multiple of
$e^{-\overline\mu r}$.
All the vectors $\mathbf1_{(a,b)}e^{-\overline\mu r}$
would consequently belong to that complement.  Their span
is dense in $L^2(0,L)$, since multiplication by this
exponential is boundedly invertible on the finite interval.
Every function in the $B_\alpha$ range would then vanish
on $(0,L)$, contrary to the no-delay conclusion of
MA-GP-003.  This proves the distinction.
\end{proof}

This argument is a precise obstruction to dropping the Gram
factor in the finite projection.  It is not an obstruction to
its existence: the bounded inverse was proved in (MA-GP.42).
Nor does it disprove a later asymptotic trace comparison with
all factors and remainders retained.

\subsection{An explicit theta Gram kernel before the Hardy identification}

\paragraph{MA-GP-005 (the actual frequency parametrization and normal equation).}
The Gram form can also be written directly from the source
PW definition, independently of the Hardy identification.
Take $F=\mathbb Q$, $S=\{\infty\}$ and a finite frequency
level $\Lambda>0$.  Let
\[
 \begin{aligned}
 \mathcal P_\Lambda^0
   &=\{u\in L^2_{\mathrm{even}}([-\Lambda,\Lambda]):
                                       \int_{-\Lambda}^{\Lambda}u=0\},\\
 \mathscr D_\Lambda
   &=\{u\in C_c^\infty((-\Lambda,\Lambda)):
                      u\text{ even},\ u(0)=0,\ \int u=0\}.
 \end{aligned}
 \tag{MA-GP.50}
\]
Use the ordinary even $L^2$ norm in this \emph{auxiliary}
source space.  It is not the weighted local norm in the draft.
For $u\in\mathscr D_\Lambda$, put $f=\mathcal F_eu$.
The two source conditions $f(0)=\int f=0$ hold.
Poisson summation gives the actual centered global theta map
\[
 (T_\Lambda u)(t)
       =2\sqrt t\sum_{n\ge1}f(nt)
       =2t^{-1/2}\sum_{n\ge1}u(n/t),\qquad t>0.
 \tag{MA-GP.51}
\]
Its target norm is
\[
 \|F\|_H^2=\int_0^\infty|F(t)|^2
           w_{\alpha,s}(\log t)\,\frac{dt}{t},
 \qquad w_{\alpha,s}(x)=e^{2\alpha|x|}(1+x^2)^{s/2}.
 \tag{MA-GP.52}
\]
Every vector (MA-GP.51) belongs to $H$, and the closure of
its range is the actual finite source PW range.
The map $T_\Lambda$ is densely defined and closable from
$\mathcal P_\Lambda^0$ into $H$.

Let $0\le\chi\in C_c^\infty(\mathbb R_{>0})$.
The localized map $M_{\sqrt\chi}T_\Lambda$ extends boundedly
to $\mathcal P_\Lambda^0$.  Its bounded positive Gram
operator is the restriction followed by zero-mean projection
of the explicit finite rational-dilation kernel
\[
 G_{\chi,\Lambda}u(y)=P_\Lambda^0\left[
 2\sum_{m,n\ge1}\frac1m\,
       \chi(m/y)w_{\alpha,s}(\log(m/y))u(ny/m)
                         \right],\qquad 0<y<\Lambda.
 \tag{MA-GP.53}
\]
Here $u$ is extended evenly by zero outside the band and
$P_\Lambda^0$ is the ordinary $L^2$ projection onto the
zero-integral even subspace.  All sums are finite for a
fixed compact output cutoff, uniformly in $y$ where they
are nonzero.

\begin{proof}
The test domain is dense in $\mathcal P_\Lambda^0$:
approximate by smooth even functions, cut away a shrinking
neighborhood of zero, and subtract their integral times
one fixed smooth even annular function of integral one.
The corrections tend to zero in $L^2$.
Poisson summation applies to $f$ and its Fourier transform
$u$, since both zeroth moments vanish, and gives (MA-GP.51).
The second expression is zero for $t<1/\Lambda$.
The first decreases faster than every power for $t\to\infty$,
uniformly after every derivative, by the Schwartz estimates
for $f$.  Thus it belongs to every target (MA-GP.52).

For this one-place source the absolute and invariant PW
support conditions coincide with the even Fourier band.
Definitions~7.3.2.1 and~7.3.2.3, pp.~351 and~355,
and Definition~11.2.1.3, pp.~488--489, give precisely the
range closure stated here.  Allowing smooth transforms
whose support meets the endpoints of the closed band
does not change it: a dilation into a slightly smaller
band converges in Schwartz topology and preserves the
two zero conditions.

If $\operatorname{supp}\chi\subset[a,b]$, the second
sum in (MA-GP.51) has at most $\lfloor\Lambda b\rfloor$
terms.  Cauchy--Schwarz and the change of variable $y=n/t$
give a finite bound by $\|u\|_2^2$ for
$\|M_{\sqrt\chi}T_\Lambda u\|_H^2$; the weight is bounded
on this compact interval.  Thus each output localization
is bounded on the auxiliary source.
If $u_j\to0$ in that source and $T_\Lambda u_j\to F$
in $H$, every such localization of $F$ is zero.
Hence $F=0$, proving closability.

For the Gram calculation, finite summation and the same
substitution yield
\[
 \begin{aligned}
 \langle \chi T_\Lambda u,T_\Lambda v\rangle_H
  &=4\sum_{m,n\ge1}\frac1n\int_0^\Lambda
       \chi(n/y)w_{\alpha,s}(\log(n/y))
                       u(y)\overline{v(my/n)}\,dy.
 \end{aligned}
 \tag{MA-GP.54}
\]
Changing the variable once more to the argument of $v$
and using the even source norm $2\int_0^\Lambda|u|^2$
gives (MA-GP.53).  The zero-mean projection appears because
the source is $\mathcal P_\Lambda^0$, rather than all even
band functions.  Both $m$ and $n$ are at most $\Lambda b$
in a nonzero term, proving uniform finiteness.
\end{proof}

Let $\overline T_\Lambda$ be this closed theta map and
$G_\Lambda=\overline T_\Lambda^{\,*}\overline T_\Lambda$
its positive self-adjoint Gram operator.  If
$\overline T_\Lambda=V G_\Lambda^{1/2}$ is its polar
decomposition, the actual global orthogonal PW projection is
\[
 Q_\Lambda
     =\operatorname*{s-lim}_{\varepsilon\downarrow0}
          V\,G_\Lambda(G_\Lambda+\varepsilon)^{-1}V^*.
 \tag{MA-GP.55}
\]
On the adjoint domain the bounded operator inside this limit
agrees with
$\overline T_\Lambda(G_\Lambda+\varepsilon)^{-1}
\overline T_\Lambda^{\,*}$.
Indeed polar decomposition gives this identity, with norm
at most one, and the spectral theorem gives the strong limit
$VV^*$ onto the range closure.  This is the domain-correct
version of a normal-equation inverse.  An unqualified formula
using $T(T^*T)^{-1}T^*$ would require additional hypotheses.
The localized forms (MA-GP.54) increase to the full target
norm on the core under a nonnegative compact exhaustion;
their uncut double expansion must not be interchanged with
the integral without a convergence proof.

\paragraph{MA-GP-006 (the frequency Gram is substantive, already at one place).}
The localized Gram kernel in (MA-GP.53) has genuine
off-diagonal rational-dilation terms, even on the two-moment
test domain.  Moreover, for every $\alpha>0$ and $s\ge0$,
$T_\Lambda$ is unbounded from the auxiliary ordinary
$L^2$ source in (MA-GP.50) into (MA-GP.52).
Neither assertion claims an unbounded inverse in a different
weighted local norm.

\begin{proof}
For the first assertion choose the output cutoff $\chi$
supported very close to $t_0=3/\Lambda$.
Take a positive input bump $u_+$ supported near $\Lambda/3$
and a positive bump $v_+$ near $2\Lambda/3$, aligned by
$v_+(2y)=u_+(y)$ on a smaller interval.
Correct the integral of $u_+$ by a bump near $\Lambda/6$,
and that of $v_+$ by a bump near $5\Lambda/6$.
Choose all four intervals sufficiently small and disjoint.
On the support of $\chi$ none of the correction intervals
is sampled by $n/t$ for an integer $n$; the only relevant
samples are near the multiples of $\Lambda/3$.
Thus the even zero-integral corrections $u,v$ are orthogonal
in the local source norm, while
\[
                    \langle\chi T_\Lambda u,T_\Lambda v\rangle_H>0.
 \tag{MA-GP.56}
\]
They vanish near the origin and hence belong to the test
domain.  This exhibits an actual off-diagonal Gram coupling.

For the second assertion choose an even nonnegative
$\psi\in C_c^\infty((-\Lambda,\Lambda))$, supported away
from zero, with nonzero integral.  Put $\Psi=\mathcal F_e\psi$.
Choose another fixed even annular test $\psi_0$ with
nonzero integral, and let
\[
 u_R(\xi)=\psi(\xi)\cos(2\pi R\xi)-c_R\psi_0(\xi),
 \qquad c_R=\frac{\int\psi(\xi)\cos(2\pi R\xi)\,d\xi}
                         {\int\psi_0(\xi)\,d\xi}.
 \tag{MA-GP.57}
\]
Then $u_R\in\mathscr D_\Lambda$, its source norm is
uniformly bounded, and $c_R=O_N(R^{-N})$ for every $N$.
Its Fourier transform is
\[
 f_R(t)=\tfrac12\bigl(\Psi(t-R)+\Psi(t+R)\bigr)
                                  -c_R(\mathcal F_e\psi_0)(t).
 \tag{MA-GP.58}
\]
Since $\Psi(0)>0$, choose $d>0$ with $\Psi(t)>0$
uniformly on $[-d,d]$.  For $t\in[R-d,R+d]$ and large
$R$, the $n=1$ summand of $\sum_{n\ge1}f_R(nt)$ is
bounded below by a fixed positive constant.  The remaining
summands are $O_N(R^{-N})$ after summing, by the Schwartz
decay at $nt-R\gg(n-1)R$ and at $nt+R$; the correction
has the same bound.  Consequently
\[
 \|T_\Lambda u_R\|_H^2
   \ge C\int_{R-d}^{R+d}w_{\alpha,s}(\log t)\,dt
   \ge C'R^{2\alpha}(1+\log^2R)^{s/2}.
 \tag{MA-GP.59}
\]
This diverges, proving the assertion with the specified
source norm and at the fixed frequency level $\Lambda$.
\end{proof}

These formulas isolate the actual finite-level normal equation.
The bounded half-line Gram inverse in MA-GP-004 and the
possibly unbounded auxiliary frequency Gram in MA-GP-005
are two different parametrizations of the same range, with
different source norms; they are not contradictory claims.
Neither may be silently replaced by the identity when
transporting a weighted orthogonal projection to an additive
Fourier cutoff.  The actual fixed-position condensation limit
is handled separately by MA-GP-001.  Any asymptotic local
trace comparison must retain the projection factors just
identified or prove that their contribution has the required
limit.

\paragraph{Grothendieck placement and source scope.}
The Gram inverse is an ordinary Hilbert-space normalization
of the actual range inclusion.  Its positivity comes from
the scalar norm weight $R$; it supplies no arithmetic
positive form on the full zero jets.  The finite-length
$\mathbb C[z]$-modules at those zeros retain the multiplicities
already identified in MA-WP, while a change of Hilbert metric
changes orthogonal representatives.  The actual trace-ideal
approximation in MA-GP-001 and the range-normalized projection
here fit the inherited cote~1 nuclear-trace comparison,
scans~079--080 and~111--112, through BR-TR-01/02 and
F-SM-11/12.  Those scan locators are not newly inspected here.
The finite-prime unitary factors in MA-LA-013 have their own
rank-one character interpretation when genuine Hecke coefficients
are supplied; the Gram calculation does not construct an
arithmetic correspondence, a new Artin motive, or an
$\ell$-adic weight statement.  The explicit Hardy range used
in (MA-GP.46)--(MA-GP.49) is exactly the theorem proved in
MA-GP-003, with the Burnol/Beurling primary-source scope stated
there; the normal-equation and frequency-kernel proofs above
are supplied in full.

\subsection{A sharp-band comparison and an actual regularized trace transport}

\paragraph{MA-GP-007 (the same-coordinate sharp-band candidate has a different range).}
The distinction from an additive Fourier cutoff persists even if its
frequency profile is a sharp characteristic function, so that it really
is an orthogonal projection.  This observation concerns a precisely
specified candidate identification; it is not attributed to a printed
operator equality in Masullo.

For the weight in (MA-GP.52) define the unitary radial map
\[
 (\mathcal U_w f)(t)
       =\sqrt{\frac{2t}{w_{\alpha,s}(\log t)}}\,f(t),\qquad t>0,
 \quad
 \mathcal U_w:L^2_{\mathrm{even}}(\mathbb R)\longrightarrow H.
 \tag{MA-GP.60}
\]
For any finite $L>0$ the same-coordinate additive candidate is
\[
 \Pi_L=\mathcal U_w\mathcal F_e
         M_{\mathbf1_{[-L,L]}}\mathcal F_e\mathcal U_w^{-1}.
 \tag{MA-GP.61}
\]
For every finite actual PW level $\Lambda>0$,
\[
               J_\Lambda\cap\operatorname{ran}\Pi_L=\{0\}.
 \tag{MA-GP.62}
\]
In particular $Q_\Lambda\ne\Pi_L$.  The same conclusion holds
after imposing the continuous zero-integral condition on the
auxiliary frequency band.  No choice of two finite positive
bandwidths changes this range obstruction.

\begin{proof}
The factor in (MA-GP.60) gives
$\|\mathcal U_w f\|_H^2=2\int_0^\infty|f(t)|^2dt$,
so it is unitary.  Every vector in $J_\Lambda$ is zero for
$0<t<1/\Lambda$, by (MA-GP.51) and closedness of this
support subspace in $H$.
If $F=\mathcal U_w f$ also belongs to $\operatorname{ran}\Pi_L$,
then $\widehat f$ is an $L^2$ function supported in $[-L,L]$.
It is in $L^1$ on that finite interval, and
\[
                 f(z)=\int_{-L}^{L}\widehat f(\xi)
                                         e^{2\pi iz\xi}\,d\xi
 \tag{MA-GP.63}
\]
is entire.  Since the radial factor is strictly positive for
$t>0$, $f$ vanishes almost everywhere on $(0,1/\Lambda)$.
Continuity and the identity theorem give $f=0$.
The actual range $J_\Lambda$ is nonzero, by a nonzero test
in (MA-GP.50) and the injectivity of the smooth theta map.
Thus the two projections cannot agree.  Passing to a smaller
moment-zero additive range preserves the zero intersection.
\end{proof}

The unweighted centering factor $\sqrt{2t}$ gives the same
support argument wherever both spaces are compared as actual
functions.  Formula (MA-GP.60) additionally puts the comparison
in one specified Hilbert space.  This argument rules out the
literal same-coordinate substitution (MA-GP.61); it neither
rules out a nonlocal transport nor contradicts an asymptotic
trace formula with its missing factors restored.

\paragraph{MA-GP-008 (an actual Gram-normalized local cyclic operator).}
Fix $\alpha>0$, $s\ge0$ and $\Lambda>0$, and retain the
auxiliary frequency space $\mathcal P_\Lambda^0$ and closed
theta map $\overline T=\overline T_\Lambda$ of MA-GP-005.
Write $G=\overline T^{\,*}\overline T$ and
$\overline T=VG^{1/2}$.
For $\chi,h\in C_c^\infty(\mathbb R_{>0})$ put
\[
 A=M_\chi C_h,\qquad
 (C_hF)(t)=\int_0^\infty h(r)F(t/r)\,\frac{dr}{r}.
 \tag{MA-GP.64}
\]
Thus $A\in\mathcal S_1(H)$, by the compact-kernel argument
of MA-LIFT-003.  There is a bounded trace-class operator
$K_{\chi,h}$ on $\mathcal P_\Lambda^0$ whose exact core form is
\[
       \langle A T_\Lambda u,T_\Lambda v\rangle_H
                         =\langle K_{\chi,h}u,v\rangle_{\mathcal P_\Lambda^0}.
 \tag{MA-GP.65}
\]
On positive band coordinates its explicit formula is
\[
 \begin{split}
 (K_{\chi,h}u)(y)&=P_\Lambda^0\biggl[
 2\sum_{m,n\ge1}\frac{\chi(m/y)w_{\alpha,s}(\log(m/y))}{m}\\
 &\hspace{18mm}\cdot\int_0^\infty h(r)r^{1/2}u(nry/m)\,\frac{dr}{r}
                                      \biggr].
 \end{split}
 \tag{MA-GP.66}
\]
Here $0<y<\Lambda$, and the output inside the projection is extended evenly.
If $\operatorname{supp}\chi\subset[a,b]$ and
$\operatorname{supp}h\subset[r_-,r_+]$ with $a,r_->0$,
only $m\le\Lambda b$ and $n\le\Lambda b/r_-$ contribute.

For every $\varepsilon>0$ define bounded operators
\[
 \begin{aligned}
 R_\varepsilon&=(G+\varepsilon)^{-1/2},&
 F_\varepsilon&=\overline T R_\varepsilon
                 =V\sqrt{G/(G+\varepsilon)},\\
 Q_{\Lambda,\varepsilon}&=F_\varepsilon F_\varepsilon^*,&
 L_{\Lambda,\varepsilon}(\chi,h)
       &=R_\varepsilon K_{\chi,h}R_\varepsilon
         =F_\varepsilon^*AF_\varepsilon.
 \end{aligned}
 \tag{MA-GP.67}
\]
Here $R_\varepsilon$ maps the entire auxiliary space into
$\operatorname{Dom}\overline T$, so the displayed maps
require no unverified adjoint domain.  They satisfy
\[
 \begin{gathered}
 \|F_\varepsilon\|\le1,\qquad
 \|L_{\Lambda,\varepsilon}(\chi,h)\|_1\le\|A\|_1,\\
 \operatorname{Tr}_H(AQ_{\Lambda,\varepsilon})
       =\operatorname{Tr}_{\mathcal P_\Lambda^0}
                           L_{\Lambda,\varepsilon}(\chi,h),\\
 \|AQ_{\Lambda,\varepsilon}-AQ_\Lambda\|_1
  +\|L_{\Lambda,\varepsilon}(\chi,h)-V^*AV\|_1
                                  \longrightarrow0.
 \end{gathered}
 \tag{MA-GP.68}
\]
In particular this is an actual local cyclic trace formula
for the specified global finite-level projection:
\[
 \operatorname{Tr}_H(M_\chi C_hQ_\Lambda)
    =\lim_{\varepsilon\downarrow0}
       \operatorname{Tr}_{\mathcal P_\Lambda^0}
       \bigl((G+\varepsilon)^{-1/2}
                K_{\chi,h}(G+\varepsilon)^{-1/2}\bigr).
 \tag{MA-GP.69}
\]

\begin{proof}
\emph{The explicit trace-class core form.}
In logarithmic coordinates the unweighted kernel of $A$ is
$\chi(e^x)h(e^{x-y})$, smooth with compact support in both
variables.  It is trace class.  Conjugating to the weighted
space introduces bounded multiplication factors on these
two compact coordinate projections; no differentiability of
$w_{\alpha,s}$ at zero is needed.  This proves the stated
trace-class property of $A$ directly.

For core vectors, insert (MA-GP.51) into (MA-GP.65).
The inner integral contributing to $AT_\Lambda u(t)$ is
\[
       2\chi(t)t^{-1/2}\sum_n
                   \int h(r)r^{1/2}u(nr/t)\,\frac{dr}{r}.
 \tag{MA-GP.70}
\]
Pair it with $2t^{-1/2}\sum_m v(m/t)$ in (MA-GP.52)
and put $y=m/t$.  The measure $t^{-2}dt$ becomes
$dy/m$, with the integration orientation reversed.
The factor four in the form becomes the factor two in
(MA-GP.66), because the source norm is
$2\int_0^\Lambda|u(y)|^2dy$.
The compact supports and the band bounds give precisely
the two finite summation bounds following (MA-GP.66).

For trace class, substitute $z=nry/m$ in that equation.
Before the bounded projection $P_\Lambda^0$, each term is
an integral operator with kernel, relative to $dz$,
\[
 \frac{2\,\chi(m/y)w_{\alpha,s}(\log(m/y))}
                        {\sqrt{mn}\sqrt{yz}}
                          h\!\left(\frac{mz}{ny}\right).
 \tag{MA-GP.71}
\]
On its support $y\ge1/b$ and $z\ge r_-/b$.
Insert smooth cutoffs equal to one on these portions of
$[0,\Lambda]$.  After removing the bounded factor
$\chi(m/y)w_{\alpha,s}(\log(m/y))$, the remaining kernel
extends smoothly to a neighborhood of the closed square
$[0,\Lambda]^2$.  Such a kernel is trace class: extend it
smoothly to a larger periodic square, expand in a double
Fourier series whose coefficients are absolutely summable
after repeated integration by parts, and restrict back to
the interval.  Bounded multiplication, the finite sum and
$P_\Lambda^0$ preserve this property.  This proves that
(MA-GP.66) defines a bounded $K_{\chi,h}\in\mathcal S_1$,
even though its weight factor can have a cusp.

\emph{Closure and domains.}
The original test domain is a graph core for $\overline T$
by its definition as the closure of $T_\Lambda$.
Approximating both vectors in graph norm therefore extends
(MA-GP.65) to all $u,v\in\operatorname{Dom}\overline T$:
the left side converges because $A$ is bounded and the right
side because $K_{\chi,h}$ is bounded.
In particular $A\overline T u\in\operatorname{Dom}\overline T^{\,*}$
for such $u$, and
$\overline T^{\,*}A\overline T u=K_{\chi,h}u$ there.
This adjoint-domain assertion is a consequence of the form
estimate, not an assumption in the initial calculation.

The spectral theorem gives
$\operatorname{Dom}\overline T=\operatorname{Dom}G^{1/2}$
and $\|G^{1/2}R_\varepsilon\|\le1$.
Thus $F_\varepsilon$ is a bounded contraction on its entire
source.  Substituting $R_\varepsilon u,R_\varepsilon v$
in the extended form proves the last identity of (MA-GP.67).
It proves it as an equality of bounded trace-class operators,
without writing an unrestricted product with $\overline T^{\,*}$.

In fact $\overline T$ is injective in this model.
The compact-output bounded maps of MA-GP-005 extend the
finite Poisson expression to every graph-domain vector.
If $\overline T u=0$, then
$S_u(y)=\sum_{m\ge1}u(my)=0$ almost everywhere on $y>0$.
For each $y>0$ the band support makes the sums finite, and
finite M\"obius inversion gives
\[
                     u(y)=\sum_{n\ge1}\mu(n)S_u(ny)=0
                                \quad\hbox{almost everywhere}.
 \tag{MA-GP.72}
\]
Dilations preserve null sets, so the almost-everywhere
statement survives the countably many possible indices.
Consequently $V$ is an isometry onto $J_\Lambda$.
The argument below would also work with its polar partial
isometry if a kernel were present.

\emph{Actual cyclic traces and their limit.}
Spectral calculus gives $F_\varepsilon\to V$ and
$F_\varepsilon^*\to V^*$ strongly, with uniform norm at
most one.  It also gives $Q_{\Lambda,\varepsilon}\to Q_\Lambda$
strongly; these are self-adjoint contractions.
For a fixed trace-class operator, uniformly bounded strong
convergence on the left, and adjoint strong convergence on
the right, imply trace-norm convergence.  This follows first
for rank-one operators and then by finite-rank approximation.
Apply it to $A$ to obtain both limits in (MA-GP.68) and
the uniform trace-norm bound there.
The cross-space cyclic identity for a trace-class factor
and a bounded factor gives
$\operatorname{Tr}(AF_\varepsilon F_\varepsilon^*)
 =\operatorname{Tr}(F_\varepsilon^*AF_\varepsilon)$.
Taking the proved trace-norm limits yields (MA-GP.69).
\end{proof}

\paragraph{The exact correction before taking any position or band limit.}
Let $E_\Lambda:\mathcal P_\Lambda^0\to L^2_{\mathrm{even}}(\mathbb R)$
be band inclusion and put
\[
 W_\Lambda=\mathcal U_w\mathcal F_eE_\Lambda,
 \qquad \Pi_\Lambda^0=W_\Lambda W_\Lambda^*.
 \tag{MA-GP.73}
\]
This is the sharp additive candidate with the same auxiliary
zero-integral source condition; $W_\Lambda$ is an isometry.
The actual finite comparison is the trace-class identity
\[
 \begin{split}
 \operatorname{Tr}_H\bigl(A(Q_{\Lambda,\varepsilon}-\Pi_\Lambda^0)\bigr)
   =\operatorname{Tr}_{\mathcal P_\Lambda^0}\bigl(
          R_\varepsilon K_{\chi,h}R_\varepsilon
                                  -W_\Lambda^*AW_\Lambda\bigr).
 \end{split}
 \tag{MA-GP.74}
\]
The first local term has the explicit rational-dilation kernel
(MA-GP.66), with its two normal-equation resolvents.
The second is the actual conjugated additive-band operator.
Thus this comparison retains both the theta summation and
the Gram normalization.  Neither can be removed by trace
cyclicity.  MA-GP-006 already proves that the theta map is
not an isometry for the auxiliary norm, and MA-GP-007
proves that the two finite ranges differ.
No nonzero value of the difference is asserted for every
particular test, and no asymptotic estimate for (MA-GP.74)
is inferred from the range distinction alone.

For completeness, at fixed $\chi,h,\Lambda,\varepsilon$
the ordinary Fredholm determinants satisfy
\[
 \det_{\mathcal P_\Lambda^0}(1-zL_{\Lambda,\varepsilon})
       =\det_H(1-zAQ_{\Lambda,\varepsilon}).
 \tag{MA-GP.75}
\]
This is Sylvester's identity, first for a finite-rank
approximation to $A$ and then by trace-norm continuity of
the Fredholm determinant.  The same reasoning and
(MA-GP.68) pass to the limit $\varepsilon\downarrow0$
at these fixed data, uniformly for $z$ in compact sets.
No determinant limit in an expanding position cutoff is
claimed.

The maps $F_\varepsilon$ provide a genuine Hilbert trace
transport.  They are not asserted to intertwine the source
and target representations: the weighted Gram need not
commute with the dilation action, and that action also
changes the fixed frequency level.  Likewise no continuity
of its resolvents on the polar-free LF source, or common
nuclear LF presentation, has been proved here.
The ordinary Hilbert trace-ideal and Fredholm-determinant
claims above have their own complete proofs.
The fixed-position condensation of MA-GP-001 can now be
expressed through (MA-GP.69); the remaining geometric
problem is an asymptotic estimate for the explicit
Gram-corrected cyclic operator, not the existence of such
an operator or an unverified use of $T^*$.
